% ch2 B.1 后半 (原书页 36–43 / PDF p051–p058)
% 衔接说明：本文件自上句半截处续起——原书 p050 末句 "This is called the first"
% 与本文件首词 "布里渊区。" 合成完整一句（"第一布里渊区"）。

这就是所谓的\emph{第一}布里渊区。如果愿意，我们可以采用所谓的“简约区图式”（reduced zone scheme），即把所有等能面用各种 $K$ 平移，从而把能带结构完全画在这第一区内，表示为 $k$ 的多值函数。这个区的体积显然是 $(2\pi)^3/\Omega$，因此它恰好包含 $N$ 个态，晶格的每个原胞对应一个态。

现在我们继续讨论“Jones 区”（Jones zones）或“大区”（large zones）(22) 的问题。到目前为止，我一直假定晶格势 $V(r)$ 除单纯的平移群之外没有其他对称性——我把一切都建立在简单 Bravais 格子之上。许多晶格在 Bravais 格子的每个原胞里不止有一个原子；许多晶格在 Bravais 格子之内还具有附加的转动或反演对称性，例如面心立方（f.c.c.）或体心立方（b.c.c.）。然而，正如 Jones 所证明的，这并不必然使任何 $V(K)$ 为零。可是，当晶格对称性除平直的平移和转动之外还包含螺旋轴与滑移面时，在许多情形下，$V(K)$ 会对某些平面族恒等于零。

一个极其重要的例子是金刚石格子 (23)，它具有反射-平移对称性的滑移面。找出 $V$ 的附加选择定则的一个简单而不严格的方法，是假定原子全同且球对称。众所周知，金刚石格子的 Bravais 格子是面心立方格子，具有 $\mathbf{a}_1=(a/2,\,a/2,\,0)$、$\mathbf{a}_2=(a/2,\,0,\,a/2)$ 等。相应的倒格子是一个b.c.c. 格子，其倒格矢为 $K=(2\pi/a)(1,1,1)$ 等；(200) 等；220；311；222；等等。原胞中的两个原子由矢量 $(a/4,\,a/4,\,a/4)$ 相联系。于是，在球对称假定下，它们对势的贡献恰正比于 X 射线结构因子（x-ray structure factor）
\begin{equation*}
S = 1 + e^{i(a/4,\,a/4,\,a/4)\,\cdot\,K} = 1 + \exp\left\{\frac{i\pi}{2}\,(l+m+n)\right\}
\end{equation*}
于是，当 $l+m+n=2(2N+1)$（$N$ 为整数）时，$S$ 为零：也就是说，对 200、222、420、600 等晶面为零。按照我们的微扰论计算，这意味着相应的能隙将为零，相应的布里渊区边界将会消失。

事实是，这个困难可以在更高阶的微扰论中得到补救。一般地说，即使 $V_K$ 可能为零，也会存在两个矢量 $K_1$ 与 $K_2$，使得 $K_1+K_2=K$ 且 $V_{K_1}\ne 0\ne V_{K_2}$。例如在这个例子中，$V_{111}$ 与 $V_{+1\,-1\,-1}$ 都不为零。现在，若 $k$ 靠近 $2\pi/a\,(200)$ 的平分面，一般而言 $E^0_k$ 与 $E^0_{k+2\pi/a\,(111)}$ 不会相近，于是我们可以用微扰论处理 $a_{k+K_1}$ 与 $a_{k+K_2}$。于是有
\begin{align*}
(E - E^0_{k+K})\,a_{k+K} &= \sum_{K'\ne K} V_{K'}^{*}\,a_{k+K-K'} \\
&\simeq V_{K_1}^{*}\,a_{k+K_2} + V_{K_2}^{*}\,a_{k+K_1}
\end{align*}
而
\begin{equation*}
(E^0_k - E^0_{k+K_1})\,a_{k+K_1} \simeq V_{K_1}\,a_k
\end{equation*}
所以净效果是：对 $K$ 区边界得到一个如下方程所示的\emph{有效}势：
\begin{align*}
(E - E^0_{k+K})\,a_{k+K} &\simeq V_{K}^{\mathrm{eff}}\,a_k \\
V_{K}^{\mathrm{eff}} &= \sum_{K_1+K_2=K} V_{K_1}\,\frac{1}{E^0_k - E^0_{k+K_1}}\,V_{K_2}^{*}
\end{align*}
而这个有效势恰恰会以一个真实矩阵元 $V_K$ 所会有的方式造成一个能隙。

这样我们就得到一般结果：在每个布里渊区边界处——除了某些特殊对称点，对其的讨论完全超出这些讲义的范围——都存在一个带隙。当然，这并不意味着作为能量的函数总有一个间隙，因为区边界不同部分的能隙肯定出现在不同的能量处，所以各能带一般会交叠。

在势微弱但不为零的极限下，我们可以重新得到对能带结构的一种非常简单而且常常有用的近似：即略去区边界附近的那部分能带，把自由电子能量 $\hbar^2 k^2/2m$ 看作一种能带结构。也就是说，我们取自由电子抛物面，把它折回简约区图式。

当我们把自由电子抛物线看作一种能带结构时，在二维或三维情形下，把自由电子抛物面的每一片折入第一区的过程中会发生相当复杂的事情；会得到复杂的简并、平坦之处等等。出于一般兴趣，我在这里复制了（图 5）f.c.c. 空间格子沿从原点到区边界六边形面的直线、即沿 $k=(x,x,x)$ 方向的能量对 $k$ 的曲线。如读者所见，简约区中相当复杂的能带结构，在铺展到“扩展区”（extended zone）之中时，可能离一条简单的自由电子抛物线并不远。研究那些反射到区内特殊对称点上的自由电子函数——特别是原点 $\Gamma$——使我们能够仅凭对称性猜测：这样一个点上的哪些波函数可能保持简并，而哪些只是自由电子近似的偶然巧合。例如，在 A 点，所有函数都具有 $e^{\pm i(2\pi x/a)}e^{\pm i(2\pi y/a)}e^{\pm i(2\pi z/a)}$ 的形式，并且可以把它们分离成四组，即四套“对称化平面波”（symmetrized plane waves；见图 5），列于表 2 中。这些对称化平面波在能带计算中是重要的。

%FIG{5}: 图 5　f.c.c. 空间格子沿从原点到区边界六边形面的直线（即 $k=(x,x,x)$ 方向）的能量对 $k$ 曲线（自由电子抛物面折回简约区后的结果）

\begin{center}
表 2

\begin{tabular}{clc}
\toprule
简并度 & 对称化平面波（S.P.W.） & 记号 \\
\midrule
$1$ & $\cos\dfrac{2\pi x}{a}\ \ \cos\dfrac{2\pi y}{a}\ \ \cos\dfrac{2\pi z}{a}$ & $\Gamma_1$ \\[4pt]
$1$ & $\sin\dfrac{2\pi x}{a}\ \ \sin\dfrac{2\pi y}{a}\ \ \sin\dfrac{2\pi z}{a}$ & $\Gamma_2$ \\[4pt]
$3$ & $\cos\dfrac{2\pi x}{a}\ \ \cos\dfrac{2\pi y}{a}\ \ \sin\dfrac{2\pi z}{a}$，等等 & $\Gamma_{15}$ \\[4pt]
$3$ & $\cos\dfrac{2\pi x}{a}\ \ \sin\dfrac{2\pi y}{a}\ \ \sin\dfrac{2\pi z}{a}$，等等 & $\Gamma'_{25}$ \\
\bottomrule
\end{tabular}
\end{center}

现在回到 Jones 区与布里渊区对比的问题。需要注意的一点是：二阶区边界处的能隙 $2V_K^{\mathrm{eff}}$ 很可能比一阶边界处的能隙小得多。因此，虽然我们一般预期每个边界处都有能隙，但在对称性复杂的情形下，特别大的能隙很可能只出现在真正的一阶边界处。这些边界就称为“Jones 区”的边界。以每原子态数来衡量的 Jones 区体积，与布里渊区的体积之间没有必然的联系；Jones 区只须每原子每自旋拥有某个有理数 $r>1$ 个态。例如，再回到众所周知的金刚石格子情形：第一布里渊区是一个胞，其区边界主要对应于 (111) $K$ 矢量，但该八面体的六个角尖被 (200) 边界切去。而在金刚石的 Jones 区中，八面体的角尖\emph{并未}被切去，这个胞每原子拥有 $9/8$ 个电子。

对特别重要的区边界的搜寻，可以超出寻找 $V\ne 0$ 的边界，进而扩展到那些 $V$ 特别大的边界。一般地，我们预期这样的边界乃是具有大 X 射线结构因子的边界，而后者只需察看 X 射线衍射花样便可简单地猜出（当然只当所有原子全同或近于全同时才如此，因为 X 射线与低能平面电子波所感受到的有效势颇为不同）。例如，金刚石中 (111) 的结构因子为
\begin{equation*}
F = 1 + \exp\left(\frac{3i}{2}\,\pi\right) = 1 - i,\qquad |F|^2 = 2
\end{equation*}
而对 $l+m+n=4N$（220、440、400 等）的那些晶面则为 $2$，$|F|^2=4$。这样看来，(220) 这一组晶面很可能围出一个特别重要的 Jones 区。

如果我们考察这一组倒格子矢量，就会认识到它们是一个棱长为 $4\times 2\pi/a$ 的面心立方倒格子的矢量；因此它也就是一个立方棱长为 $a/2$ 的体心格子的倒格子，原 f.c.c. 格子的每一个立方体对应于它的八个立方体。于是，由于b.c.c. 每个立方体只有 2 个原子而不是 4 个，这个格子的第一布里渊区对每个 f.c.c. 格子原子有 $8/2$ 即 4 个态，对每个金刚石格子原子有 2 个态——计入自旋总共 4 个。

看待这个准体心（pseudo-body-centered）格子的另一种方式，是把金刚石格子画出来，如图 6。

%FIG{6}: 图 6　金刚石格子。○＝f.c.c. 格子的第一平面；□＝f.c.c. 格子的第二平面；斜纹圆＝金刚石格子的中间平面

给出这些强反射的b.c.c. 格子可以这样得到：把每个原子周围四面体所缺的位置补齐，使之构成立方体，如图 7 中的圆圈 O 所示。

%FIG{7}: 图 7　补全每个原子周围四面体的空缺位置（圆圈 O）以构成立方体，即给出强反射的b.c.c. 格子

这表明恰有两倍之多的态能装进它的布里渊区；在金刚石格子的强 Jones 区中，每个原子有四个电子态。

这一切都已被详细推演过，以此作为下述量子化学规则之半经验有效性的一个极好例证——正是这条规则赋予了 Jones 区重要性：晶体倾向于选择这样的结构，使得一个强的、对称的 Jones 区能够恰好或几乎恰好被填满。其理由很容易看出：如果晶体价带中的电子数使得在自由电子模型下费米面恰好会到达 Jones 区边界，那么该边界的形成就会降低区边界附近那群自由电子的能量——办法是使被占据的电子态相对于未被占据的态能量降低。

为了定量估计这一效应，设相应的 $V_K$ 的大小为 $V$。平均而言，粗略地说，能量在费米面附近 $V$ 范围内的每个电子，其能量都被降低了 $V/2$，于是降低量为 $(V^2/2)$ 乘以费米面处的态密度。在自由电子模型下，这个态密度由下式给出：
\begin{align*}
\text{每个电子的态密度} &= \frac{dn}{dE}\,\frac{1}{n} \\
&= \frac{dk}{dE}\left(\frac{1}{n}\,\frac{dn}{dk}\right) = \frac{dk}{dE}\,\frac{3}{k_F} \\
&= \frac{3}{2E_F}
\end{align*}
于是每原子的结合能将为
\begin{equation*}
\sim \frac{4\times 3}{2\times 2E_F}\,V^2 = \frac{3V^2}{E_F}
\end{equation*}
必须提醒读者，这一估计中的数值因子实际上不可能算得比数量级更准；事实上，有关这一情形的能量学目前在相当程度上仍有争议 (24)。

这一效应可能相当大，而且很可能正是它决定着晶体结构——只要在一些真实物质中估计一下它的影响便可以看出。例如，金刚石的自由电子费米能将是 30 ev（约 2 Rydberg），而在密度小得多的晶体 Si 和 Ge 中约为其一半。（这只须由 $n=k_F^3/3\pi^2$ 算出 $k_F$ 即可得到；结果金刚石的 $k_F$ 约为 $3\times 10^{8}\,\mathrm{cm}^{-1}$，Si 约为 $2\times 10^{8}/\mathrm{cm}$。）自由电子所感受到的势的各分量的实际数值，可从 Herman (23) 与 Phillips (25) 的计算中获得；它们给出：金刚石中 $V_{220}$ 约为 0.4 ryd，即 5 ev；Si 中约为 0.15--0.2 ryd，即 2 ev。于是内聚能的增益非常大：金刚石中为 3 伏特（译注：原文单位为 volts，此处即指电子伏），约占总内聚能的一半；Si 中为 1.5 伏特。诚然这是一个非常粗略的估计，但它表明，在决定一种物质将采取何种晶体结构的、通常相当微妙的能量平衡之中，这可能是一种真实而起支配作用的效应。

一个有趣的事实是：$K=111$ 的 $V_K$ 也在金刚石和硅的带隙附近扮演重要角色，它在某些 $V_{220}$ 无法起作用的对称点——尤其是 $\Gamma$ 点——处对打开带隙至关重要。因此，对结合能的一个可观贡献也来自 $V_{111}$。这大概是几个因素的组合决定了所观察到的晶体结构的一个典型例子。

在转向更定量的内容之前，我想从这一粗略计算中引出若干物理与化学上的教益。首先，注意我们在这里为开放型、价型（valence-type）结构看到了一个物理上的理由，它完全摆脱了把定向键设想为从原子沿确定方向伸出的若干小棍棒的观念。在一个多价物质中，开放结构完全可以只是一种手段，把重要的 Jones 区边界带进费米面所处的区域。不过我应当提醒读者：作为一件化学事实，键的方向性看来确实存在；而且，一个给定物质究竟会采取具有大致正确的密度与 Jones 区形状的哪一个具体结构，其计算很可能是 $V$ 的极多傅里叶分量之间复杂相互作用的问题。眼下这通常最好留给行之有效的半经验方法与化学直觉去处理；但我们至少看到了其中涉及的某些原理，而且这些原理并不像 Pauling 可能会提示的那样，要求抛弃能带理论这套非常成功的工具。

其次，注意在这两种情形中，$V$ 实际上都小到使微扰论并不太差的程度。对金刚石来说这是相当令人吃惊的，但却是事实——金刚石巨大的能隙只是碳的整个能量标度很宽这一事实的后果，并不代表对近自由电子图像的真正大的偏离。

最后，我应当提到，这个对金刚石和硅的应用实际上是这类应用中最新的一个；历史悠久得多、研究也彻底得多的领域，是把它应用于某些合金结构——或者更确切地说，金属间化合物结构——特别是著名的 Hume-Rothery 合金。欲知较好的论述，可参见 Jones (26)。这些化合物具有各种多少有些复杂的立方结构，往往有一个对称性相当高的突出 Jones 区位于非常靠近费米面之处。它们无疑是由同样的效应引起的，但毫无疑问，其能量贡献比价结构中的要小得多。我们看到这一点的依据是：这些化合物仍保持其金属性，表明区边界并未打开多大的能隙。还值得强调的是，即使在例如金刚石的情形中，某些特殊的对称方向——即 Jones 区的角点，按对称性它们等价于简约区的中心 $\Gamma$——在近自由图像中也无法仅由 $V_{220}$ 这一个分量使之劈裂，因此 $V$ 的其他分量，特别是我们前面提到的 $V_{111}$，在打开造成绝缘行为的实际能隙方面起着支配作用。
